% --- Archon physics formalization source begin ---
% archon:physics
% archon:covers PhyXMiniProblems/problem_phyx_mini_0967.lean
% archon:source-report reports/phyx_mini/problem_phyx_mini_0967.source.json

\chapter{Physics problem phyx\_mini\_0967}
\label{ch:PhyXMiniProblems_problem_phyx_mini_0967}

\paragraph{Problem source.}
\#\# Physical scenario

A cylindrical shell with radius \textbackslash{}( R\_1 \textbackslash{}) and height \textbackslash{}( H \textbackslash{}) has charge \textbackslash{}( Q\_1 \textbackslash{}) and rotates around its axis with angular speed \textbackslash{}( \textbackslash{}omega\_1 \textbackslash{}). Inside the cylinder, far from its edges, sits a very small disk with radius \textbackslash{}( R\_2 \textbackslash{}), mass \textbackslash{}( M \textbackslash{}), and charge \textbackslash{}( Q\_2 \textbackslash{}) mounted on a pivot, spinning with a large angular velocity \textbackslash{}( \textbackslash{}vec\{\textbackslash{}omega\}\_2 \textbackslash{}) and oriented at angle \textbackslash{}( \textbackslash{}theta \textbackslash{}) with respect to the axis of the cylinder. The center of the disk is on the axis of the cylinder. The magnetic interaction between the cylinder and the disk causes a precession of the axis of the disk.

\#\# Question

The magnetic moment of the disk has magnitude \textbackslash{}( \textbackslash{}mu = \textbackslash{}frac\{1\}\{4\} Q\_2 \textbackslash{}omega\_2 R\_2\textasciicircum{}2 \textbackslash{}). What is the magnitude of the torque exerted on the disk?

\#\# Answer choices

- A: A: \textbackslash{}frac\{\textbackslash{}mu\_0 Q\_1 Q\_2 \textbackslash{}omega\_1 \textbackslash{}omega\_2 R\_2 \textbackslash{}sin \textbackslash{}theta\}\{8 \textbackslash{}pi H\}

- B: B: \textbackslash{}frac\{\textbackslash{}mu\_0 Q\_1 Q\_2 \textbackslash{}omega\_1 \textbackslash{}omega\_2 R\_2\textasciicircum{}2 \textbackslash{}sin \textbackslash{}theta\}\{8 \textbackslash{}pi H\}

- C: C: \textbackslash{}frac\{\textbackslash{}mu\_0 Q\_1 Q\_2 \textbackslash{}omega\_1 \textbackslash{}omega\_2 R\_2\textasciicircum{}2 \textbackslash{}cos \textbackslash{}theta\}\{8 \textbackslash{}pi H\}

- D: D: \textbackslash{}frac\{\textbackslash{}mu\_0 Q\_1 Q\_2 \textbackslash{}omega\_1 \textbackslash{}omega\_2 R\_2\textasciicircum{}2\}\{8 \textbackslash{}pi H\}

\#\# Figure caption (auxiliary; use the image as primary evidence)

The image is a diagram of a cylindrical object with various elements labeled, likely from a physics or mechanics context. Here's a detailed description:

1. **Cylinder:**
   - The cylinder has a height labeled \textbackslash{}(H\textbackslash{}).
   - The radius of the cylinder is labeled \textbackslash{}(R\_1\textbackslash{}).
   - The central vertical line is labeled as the "Axis of cylinder."

2. **Internal Elements:**
   - Inside the cylinder, there's an inner circular object labeled \textbackslash{}(Q\_2\textbackslash{}) with radius \textbackslash{}(R\_2\textbackslash{}).
   - The inner circle has an angle \textbackslash{}(\textbackslash{}theta\textbackslash{}) and an angular velocity \textbackslash{}(\textbackslash{}vec\{\textbackslash{}omega\}\_2\textbackslash{}).
   - There's a vector \textbackslash{}(\textbackslash{}Omega\textbackslash{}) pointing upward from the center of the inner circle.

3. **Arrows and Labels:**
   - \textbackslash{}(\textbackslash{}vec\{\textbackslash{}omega\}\_1\textbackslash{}) is an angular velocity vector pointing along the side of the cylinder.
   - There are curved arrows indicating rotation around both axes.

4. **Additional Labels:**
   - \textbackslash{}(Q\_1\textbackslash{}) is labeled at the bottom of the cylinder.
   - \textbackslash{}(M\textbackslash{}) appears to be labeled near the axis or as part of the circle inside the cylinder.

5. **Other Features:**
   - The diagram includes arrows and circular paths indicating motion and relationships between different components within and outside the cylinder.

This diagram likely represents a mechanical system involving rotational motion, possibly in the context of angular momentum or dynamics.

\#\# Dataset metadata

Magnetism; Physical Model Grounding Reasoning, Multi-Formula Reasoning

\paragraph{Recorded answer/context.}
B: B: \textbackslash{}frac\{\textbackslash{}mu\_0 Q\_1 Q\_2 \textbackslash{}omega\_1 \textbackslash{}omega\_2 R\_2\textasciicircum{}2 \textbackslash{}sin \textbackslash{}theta\}\{8 \textbackslash{}pi H\}

\paragraph{Figure/image path.}
/root/proposal\_for\_physic/science-mango/phyx\_mini\_run/phyx\_data/test\_image/967.png

\paragraph{Formalization target.}
create a compiling Lean file with sorry bodies at `PhyXMiniProblems/problem\_phyx\_mini\_0967.lean`. The Lean declarations must preserve the physical quantities, dimensions or dimensional roles, figure labels, governing-law hypotheses, and final relation expressed by this problem.
Use Mathlib/Physlib names found through LeanExplore where available. If a physics API is missing, introduce faithful local abstractions rather than scalar placeholder aliases.

\begin{theorem}[Physics formalization target]
\label{thm:physics:phyx_mini_0967:target}
\lean{PhyXMiniProblems.ProblemPhyXMini0967.problem_phyx_mini_0967}
\uses{def:physics:phyx-mini-0967:phyxminiproblems-problemphyxmini0967-lengthinmeters, def:physics:phyx-mini-0967:phyxminiproblems-problemphyxmini0967-chargemagnitudeincoulombs, def:physics:phyx-mini-0967:phyxminiproblems-problemphyxmini0967-angularspeedinradianspersecond, def:physics:phyx-mini-0967:phyxminiproblems-problemphyxmini0967-vacuumpermeabilityinnewtonsperamperesquared, def:physics:phyx-mini-0967:phyxminiproblems-problemphyxmini0967-torquemagnitudeinnewtonmeters, def:physics:phyx-mini-0967:phyxminiproblems-problemphyxmini0967-rotatingcylinderdisksetup, def:physics:phyx-mini-0967:phyxminiproblems-problemphyxmini0967-matchesrotatingcylinderdiskscenario, def:physics:phyx-mini-0967:phyxminiproblems-problemphyxmini0967-matchessuppliedrotatingcylinderdiskfigure, def:physics:phyx-mini-0967:phyxminiproblems-problemphyxmini0967-hasphysicalrotatingcylinderdiskparameters, def:physics:phyx-mini-0967:phyxminiproblems-problemphyxmini0967-matchescylinderdiskorientationgeometry, def:physics:phyx-mini-0967:phyxminiproblems-problemphyxmini0967-satisfiesrotatingchargedcylinderfieldlaws, def:physics:phyx-mini-0967:phyxminiproblems-problemphyxmini0967-satisfiesfinitecylinderfieldapproximation, def:physics:phyx-mini-0967:phyxminiproblems-problemphyxmini0967-satisfiesspinningchargeddiskmagneticmomentlaw, def:physics:phyx-mini-0967:phyxminiproblems-problemphyxmini0967-satisfiesmagneticdipoletorquelaw, def:physics:phyx-mini-0967:phyxminiproblems-problemphyxmini0967-satisfiesfiniteapparatustorqueapproximation, def:physics:phyx-mini-0967:phyxminiproblems-problemphyxmini0967-displayedtorquemagnitudeinnewtonmeters, def:physics:phyx-mini-0967:phyxminiproblems-problemphyxmini0967-recordeddatasetanswer}
The assigned autoformalize agent should translate this physics problem into Lean declarations in the covered file, with theorem and lemma proof bodies written as `by sorry`.
\end{theorem}
\begin{proof}
This is an autoformalization task, not a proof task. Produce statements that can later be proved by the physics prover without weakening the physical model.
\end{proof}

% --- Archon declaration topology begin ---
\section{Declaration topology}
\begin{definition}[Spatial Vector]
  \label{def:physics:phyx-mini-0967:phyxminiproblems-problemphyxmini0967-spatialvector}
  \lean{PhyXMiniProblems.ProblemPhyXMini0967.SpatialVector}
  Three-dimensional Euclidean vectors used for coherent-SI readouts
\end{definition}
\begin{proof}
  This is the declared model definition of Spatial Vector.
\end{proof}
\begin{definition}[angular Speed Dimension]
  \label{def:physics:phyx-mini-0967:phyxminiproblems-problemphyxmini0967-angularspeeddimension}
  \lean{PhyXMiniProblems.ProblemPhyXMini0967.angularSpeedDimension}
  Angular speed has inverse-time dimension; radians are dimensionless
\end{definition}
\begin{proof}
  This is the declared model definition of angular Speed Dimension.
\end{proof}
\begin{definition}[surface Current Density Dimension]
  \label{def:physics:phyx-mini-0967:phyxminiproblems-problemphyxmini0967-surfacecurrentdensitydimension}
  \lean{PhyXMiniProblems.ProblemPhyXMini0967.surfaceCurrentDensityDimension}
  Azimuthal surface current per axial length has dimension C T⁻¹ L⁻¹
\end{definition}
\begin{proof}
  This is the declared model definition of surface Current Density Dimension.
\end{proof}
\begin{definition}[vacuum Permeability Dimension]
  \label{def:physics:phyx-mini-0967:phyxminiproblems-problemphyxmini0967-vacuumpermeabilitydimension}
  \lean{PhyXMiniProblems.ProblemPhyXMini0967.vacuumPermeabilityDimension}
  Vacuum permeability has SI dimension M L C⁻²
\end{definition}
\begin{proof}
  This is the declared model definition of vacuum Permeability Dimension.
\end{proof}
\begin{definition}[magnetic Flux Density Dimension]
  \label{def:physics:phyx-mini-0967:phyxminiproblems-problemphyxmini0967-magneticfluxdensitydimension}
  \lean{PhyXMiniProblems.ProblemPhyXMini0967.magneticFluxDensityDimension}
  Magnetic flux density (tesla) has dimension M T⁻¹ C⁻¹
\end{definition}
\begin{proof}
  This is the declared model definition of magnetic Flux Density Dimension.
\end{proof}
\begin{definition}[magnetic Moment Dimension]
  \label{def:physics:phyx-mini-0967:phyxminiproblems-problemphyxmini0967-magneticmomentdimension}
  \lean{PhyXMiniProblems.ProblemPhyXMini0967.magneticMomentDimension}
  Magnetic dipole moment has dimension C L² T⁻¹
\end{definition}
\begin{proof}
  This is the declared model definition of magnetic Moment Dimension.
\end{proof}
\begin{definition}[torque Dimension]
  \label{def:physics:phyx-mini-0967:phyxminiproblems-problemphyxmini0967-torquedimension}
  \lean{PhyXMiniProblems.ProblemPhyXMini0967.torqueDimension}
  Torque has the energy dimension M L² T⁻²
\end{definition}
\begin{proof}
  This is the declared model definition of torque Dimension.
\end{proof}
\begin{definition}[Mass Magnitude]
  \label{def:physics:phyx-mini-0967:phyxminiproblems-problemphyxmini0967-massmagnitude}
  \lean{PhyXMiniProblems.ProblemPhyXMini0967.MassMagnitude}
  A nonnegative, unit-independent physical mass
\end{definition}
\begin{proof}
  This is the declared model definition of Mass Magnitude.
\end{proof}
\begin{definition}[Length Magnitude]
  \label{def:physics:phyx-mini-0967:phyxminiproblems-problemphyxmini0967-lengthmagnitude}
  \lean{PhyXMiniProblems.ProblemPhyXMini0967.LengthMagnitude}
  A nonnegative, unit-independent physical length
\end{definition}
\begin{proof}
  This is the declared model definition of Length Magnitude.
\end{proof}
\begin{definition}[Charge Magnitude]
  \label{def:physics:phyx-mini-0967:phyxminiproblems-problemphyxmini0967-chargemagnitude}
  \lean{PhyXMiniProblems.ProblemPhyXMini0967.ChargeMagnitude}
  A nonnegative electric-charge magnitude
\end{definition}
\begin{proof}
  This is the declared model definition of Charge Magnitude.
\end{proof}
\begin{definition}[Angular Speed Magnitude]
  \label{def:physics:phyx-mini-0967:phyxminiproblems-problemphyxmini0967-angularspeedmagnitude}
  \lean{PhyXMiniProblems.ProblemPhyXMini0967.AngularSpeedMagnitude}
  \uses{def:physics:phyx-mini-0967:phyxminiproblems-problemphyxmini0967-angularspeeddimension}
  A nonnegative angular-speed magnitude
\end{definition}
\begin{proof}
  This is the declared model definition of Angular Speed Magnitude.
\end{proof}
\begin{definition}[Surface Current Density Magnitude]
  \label{def:physics:phyx-mini-0967:phyxminiproblems-problemphyxmini0967-surfacecurrentdensitymagnitude}
  \lean{PhyXMiniProblems.ProblemPhyXMini0967.SurfaceCurrentDensityMagnitude}
  \uses{def:physics:phyx-mini-0967:phyxminiproblems-problemphyxmini0967-surfacecurrentdensitydimension}
  A nonnegative azimuthal surface-current density magnitude
\end{definition}
\begin{proof}
  This is the declared model definition of Surface Current Density Magnitude.
\end{proof}
\begin{definition}[Vacuum Permeability Magnitude]
  \label{def:physics:phyx-mini-0967:phyxminiproblems-problemphyxmini0967-vacuumpermeabilitymagnitude}
  \lean{PhyXMiniProblems.ProblemPhyXMini0967.VacuumPermeabilityMagnitude}
  \uses{def:physics:phyx-mini-0967:phyxminiproblems-problemphyxmini0967-vacuumpermeabilitydimension}
  A nonnegative vacuum-permeability magnitude
\end{definition}
\begin{proof}
  This is the declared model definition of Vacuum Permeability Magnitude.
\end{proof}
\begin{definition}[Magnetic Flux Density Magnitude]
  \label{def:physics:phyx-mini-0967:phyxminiproblems-problemphyxmini0967-magneticfluxdensitymagnitude}
  \lean{PhyXMiniProblems.ProblemPhyXMini0967.MagneticFluxDensityMagnitude}
  \uses{def:physics:phyx-mini-0967:phyxminiproblems-problemphyxmini0967-magneticfluxdensitydimension}
  A nonnegative magnetic-flux-density magnitude
\end{definition}
\begin{proof}
  This is the declared model definition of Magnetic Flux Density Magnitude.
\end{proof}
\begin{definition}[Magnetic Flux Density Vector]
  \label{def:physics:phyx-mini-0967:phyxminiproblems-problemphyxmini0967-magneticfluxdensityvector}
  \lean{PhyXMiniProblems.ProblemPhyXMini0967.MagneticFluxDensityVector}
  \uses{def:physics:phyx-mini-0967:phyxminiproblems-problemphyxmini0967-spatialvector, def:physics:phyx-mini-0967:phyxminiproblems-problemphyxmini0967-magneticfluxdensitydimension}
  A unit-independent magnetic-flux-density vector
\end{definition}
\begin{proof}
  This is the declared model definition of Magnetic Flux Density Vector.
\end{proof}
\begin{definition}[Magnetic Moment Magnitude]
  \label{def:physics:phyx-mini-0967:phyxminiproblems-problemphyxmini0967-magneticmomentmagnitude}
  \lean{PhyXMiniProblems.ProblemPhyXMini0967.MagneticMomentMagnitude}
  \uses{def:physics:phyx-mini-0967:phyxminiproblems-problemphyxmini0967-magneticmomentdimension}
  A nonnegative magnetic-dipole-moment magnitude
\end{definition}
\begin{proof}
  This is the declared model definition of Magnetic Moment Magnitude.
\end{proof}
\begin{definition}[Magnetic Moment Vector]
  \label{def:physics:phyx-mini-0967:phyxminiproblems-problemphyxmini0967-magneticmomentvector}
  \lean{PhyXMiniProblems.ProblemPhyXMini0967.MagneticMomentVector}
  \uses{def:physics:phyx-mini-0967:phyxminiproblems-problemphyxmini0967-spatialvector, def:physics:phyx-mini-0967:phyxminiproblems-problemphyxmini0967-magneticmomentdimension}
  A unit-independent magnetic-dipole-moment vector
\end{definition}
\begin{proof}
  This is the declared model definition of Magnetic Moment Vector.
\end{proof}
\begin{definition}[Torque Magnitude]
  \label{def:physics:phyx-mini-0967:phyxminiproblems-problemphyxmini0967-torquemagnitude}
  \lean{PhyXMiniProblems.ProblemPhyXMini0967.TorqueMagnitude}
  \uses{def:physics:phyx-mini-0967:phyxminiproblems-problemphyxmini0967-torquedimension}
  A nonnegative magnetic-torque magnitude
\end{definition}
\begin{proof}
  This is the declared model definition of Torque Magnitude.
\end{proof}
\begin{definition}[Torque Vector]
  \label{def:physics:phyx-mini-0967:phyxminiproblems-problemphyxmini0967-torquevector}
  \lean{PhyXMiniProblems.ProblemPhyXMini0967.TorqueVector}
  \uses{def:physics:phyx-mini-0967:phyxminiproblems-problemphyxmini0967-spatialvector, def:physics:phyx-mini-0967:phyxminiproblems-problemphyxmini0967-torquedimension}
  A unit-independent magnetic-torque vector
\end{definition}
\begin{proof}
  This is the declared model definition of Torque Vector.
\end{proof}
\begin{definition}[nonnegative SIReadout]
  \label{def:physics:phyx-mini-0967:phyxminiproblems-problemphyxmini0967-nonnegativesireadout}
  \lean{PhyXMiniProblems.ProblemPhyXMini0967.nonnegativeSIReadout}
  Read a nonnegative dimensionful scalar in coherent SI units
\end{definition}
\begin{proof}
  This is the declared model definition of nonnegative SIReadout.
\end{proof}
\begin{definition}[vector SIReadout]
  \label{def:physics:phyx-mini-0967:phyxminiproblems-problemphyxmini0967-vectorsireadout}
  \lean{PhyXMiniProblems.ProblemPhyXMini0967.vectorSIReadout}
  \uses{def:physics:phyx-mini-0967:phyxminiproblems-problemphyxmini0967-spatialvector}
  Read a dimensionful spatial vector in coherent SI units
\end{definition}
\begin{proof}
  This is the declared model definition of vector SIReadout.
\end{proof}
\begin{definition}[mass In Kilograms]
  \label{def:physics:phyx-mini-0967:phyxminiproblems-problemphyxmini0967-massinkilograms}
  \lean{PhyXMiniProblems.ProblemPhyXMini0967.massInKilograms}
  \uses{def:physics:phyx-mini-0967:phyxminiproblems-problemphyxmini0967-massmagnitude, def:physics:phyx-mini-0967:phyxminiproblems-problemphyxmini0967-nonnegativesireadout}
  Read mass in kilograms
\end{definition}
\begin{proof}
  This is the declared model definition of mass In Kilograms.
\end{proof}
\begin{definition}[length In Meters]
  \label{def:physics:phyx-mini-0967:phyxminiproblems-problemphyxmini0967-lengthinmeters}
  \lean{PhyXMiniProblems.ProblemPhyXMini0967.lengthInMeters}
  \uses{def:physics:phyx-mini-0967:phyxminiproblems-problemphyxmini0967-lengthmagnitude, def:physics:phyx-mini-0967:phyxminiproblems-problemphyxmini0967-nonnegativesireadout}
  Read length in metres
\end{definition}
\begin{proof}
  This is the declared model definition of length In Meters.
\end{proof}
\begin{definition}[charge Magnitude In Coulombs]
  \label{def:physics:phyx-mini-0967:phyxminiproblems-problemphyxmini0967-chargemagnitudeincoulombs}
  \lean{PhyXMiniProblems.ProblemPhyXMini0967.chargeMagnitudeInCoulombs}
  \uses{def:physics:phyx-mini-0967:phyxminiproblems-problemphyxmini0967-chargemagnitude, def:physics:phyx-mini-0967:phyxminiproblems-problemphyxmini0967-nonnegativesireadout}
  Read charge magnitude in coulombs
\end{definition}
\begin{proof}
  This is the declared model definition of charge Magnitude In Coulombs.
\end{proof}
\begin{definition}[angular Speed In Radians Per Second]
  \label{def:physics:phyx-mini-0967:phyxminiproblems-problemphyxmini0967-angularspeedinradianspersecond}
  \lean{PhyXMiniProblems.ProblemPhyXMini0967.angularSpeedInRadiansPerSecond}
  \uses{def:physics:phyx-mini-0967:phyxminiproblems-problemphyxmini0967-angularspeedmagnitude, def:physics:phyx-mini-0967:phyxminiproblems-problemphyxmini0967-nonnegativesireadout}
  Read angular-speed magnitude in radians per second
\end{definition}
\begin{proof}
  This is the declared model definition of angular Speed In Radians Per Second.
\end{proof}
\begin{definition}[surface Current Density In Amperes Per Meter]
  \label{def:physics:phyx-mini-0967:phyxminiproblems-problemphyxmini0967-surfacecurrentdensityinamperespermeter}
  \lean{PhyXMiniProblems.ProblemPhyXMini0967.surfaceCurrentDensityInAmperesPerMeter}
  \uses{def:physics:phyx-mini-0967:phyxminiproblems-problemphyxmini0967-surfacecurrentdensitymagnitude, def:physics:phyx-mini-0967:phyxminiproblems-problemphyxmini0967-nonnegativesireadout}
  Read azimuthal surface-current density in amperes per metre
\end{definition}
\begin{proof}
  This is the declared model definition of surface Current Density In Amperes Per Meter.
\end{proof}
\begin{definition}[vacuum Permeability In Newtons Per Ampere Squared]
  \label{def:physics:phyx-mini-0967:phyxminiproblems-problemphyxmini0967-vacuumpermeabilityinnewtonsperamperesquared}
  \lean{PhyXMiniProblems.ProblemPhyXMini0967.vacuumPermeabilityInNewtonsPerAmpereSquared}
  \uses{def:physics:phyx-mini-0967:phyxminiproblems-problemphyxmini0967-vacuumpermeabilitymagnitude, def:physics:phyx-mini-0967:phyxminiproblems-problemphyxmini0967-nonnegativesireadout}
  Read vacuum permeability in newtons per ampere squared
\end{definition}
\begin{proof}
  This is the declared model definition of vacuum Permeability In Newtons Per Ampere Squared.
\end{proof}
\begin{definition}[magnetic Flux Density In Teslas]
  \label{def:physics:phyx-mini-0967:phyxminiproblems-problemphyxmini0967-magneticfluxdensityinteslas}
  \lean{PhyXMiniProblems.ProblemPhyXMini0967.magneticFluxDensityInTeslas}
  \uses{def:physics:phyx-mini-0967:phyxminiproblems-problemphyxmini0967-magneticfluxdensitymagnitude, def:physics:phyx-mini-0967:phyxminiproblems-problemphyxmini0967-nonnegativesireadout}
  Read magnetic-flux-density magnitude in teslas
\end{definition}
\begin{proof}
  This is the declared model definition of magnetic Flux Density In Teslas.
\end{proof}
\begin{definition}[magnetic Flux Density Vector In Teslas]
  \label{def:physics:phyx-mini-0967:phyxminiproblems-problemphyxmini0967-magneticfluxdensityvectorinteslas}
  \lean{PhyXMiniProblems.ProblemPhyXMini0967.magneticFluxDensityVectorInTeslas}
  \uses{def:physics:phyx-mini-0967:phyxminiproblems-problemphyxmini0967-spatialvector, def:physics:phyx-mini-0967:phyxminiproblems-problemphyxmini0967-magneticfluxdensityvector, def:physics:phyx-mini-0967:phyxminiproblems-problemphyxmini0967-vectorsireadout}
  Read a magnetic-flux-density vector in tesla coordinates
\end{definition}
\begin{proof}
  This is the declared model definition of magnetic Flux Density Vector In Teslas.
\end{proof}
\begin{definition}[magnetic Moment Magnitude In Ampere Square Meters]
  \label{def:physics:phyx-mini-0967:phyxminiproblems-problemphyxmini0967-magneticmomentmagnitudeinamperesquaremeters}
  \lean{PhyXMiniProblems.ProblemPhyXMini0967.magneticMomentMagnitudeInAmpereSquareMeters}
  \uses{def:physics:phyx-mini-0967:phyxminiproblems-problemphyxmini0967-magneticmomentmagnitude, def:physics:phyx-mini-0967:phyxminiproblems-problemphyxmini0967-nonnegativesireadout}
  Read magnetic-moment magnitude in ampere-square-metres
\end{definition}
\begin{proof}
  This is the declared model definition of magnetic Moment Magnitude In Ampere Square Meters.
\end{proof}
\begin{definition}[magnetic Moment Vector In Ampere Square Meters]
  \label{def:physics:phyx-mini-0967:phyxminiproblems-problemphyxmini0967-magneticmomentvectorinamperesquaremeters}
  \lean{PhyXMiniProblems.ProblemPhyXMini0967.magneticMomentVectorInAmpereSquareMeters}
  \uses{def:physics:phyx-mini-0967:phyxminiproblems-problemphyxmini0967-spatialvector, def:physics:phyx-mini-0967:phyxminiproblems-problemphyxmini0967-magneticmomentvector, def:physics:phyx-mini-0967:phyxminiproblems-problemphyxmini0967-vectorsireadout}
  Read a magnetic-moment vector in ampere-square-metre coordinates
\end{definition}
\begin{proof}
  This is the declared model definition of magnetic Moment Vector In Ampere Square Meters.
\end{proof}
\begin{definition}[torque Magnitude In Newton Meters]
  \label{def:physics:phyx-mini-0967:phyxminiproblems-problemphyxmini0967-torquemagnitudeinnewtonmeters}
  \lean{PhyXMiniProblems.ProblemPhyXMini0967.torqueMagnitudeInNewtonMeters}
  \uses{def:physics:phyx-mini-0967:phyxminiproblems-problemphyxmini0967-torquemagnitude, def:physics:phyx-mini-0967:phyxminiproblems-problemphyxmini0967-nonnegativesireadout}
  Read torque magnitude in newton-metres
\end{definition}
\begin{proof}
  This is the declared model definition of torque Magnitude In Newton Meters.
\end{proof}
\begin{definition}[torque Vector In Newton Meters]
  \label{def:physics:phyx-mini-0967:phyxminiproblems-problemphyxmini0967-torquevectorinnewtonmeters}
  \lean{PhyXMiniProblems.ProblemPhyXMini0967.torqueVectorInNewtonMeters}
  \uses{def:physics:phyx-mini-0967:phyxminiproblems-problemphyxmini0967-spatialvector, def:physics:phyx-mini-0967:phyxminiproblems-problemphyxmini0967-torquevector, def:physics:phyx-mini-0967:phyxminiproblems-problemphyxmini0967-vectorsireadout}
  Read a torque vector in newton-metre coordinates
\end{definition}
\begin{proof}
  This is the declared model definition of torque Vector In Newton Meters.
\end{proof}
\begin{definition}[spatial Cross]
  \label{def:physics:phyx-mini-0967:phyxminiproblems-problemphyxmini0967-spatialcross}
  \lean{PhyXMiniProblems.ProblemPhyXMini0967.spatialCross}
  \uses{def:physics:phyx-mini-0967:phyxminiproblems-problemphyxmini0967-spatialvector}
  Mathlib's ordinary three-dimensional cross product transported to EuclideanSpace ℝ (Fin 3)
\end{definition}
\begin{proof}
  This is the declared model definition of spatial Cross.
\end{proof}
\begin{definition}[Cylindrical Shell Model]
  \label{def:physics:phyx-mini-0967:phyxminiproblems-problemphyxmini0967-cylindricalshellmodel}
  \lean{PhyXMiniProblems.ProblemPhyXMini0967.CylindricalShellModel}
  The idealized charge distribution on the rotating cylinder
\end{definition}
\begin{proof}
  This is the declared model definition of Cylindrical Shell Model.
\end{proof}
\begin{definition}[Disk Model]
  \label{def:physics:phyx-mini-0967:phyxminiproblems-problemphyxmini0967-diskmodel}
  \lean{PhyXMiniProblems.ProblemPhyXMini0967.DiskModel}
  The idealized mass and charge geometry of the small rotor
\end{definition}
\begin{proof}
  This is the declared model definition of Disk Model.
\end{proof}
\begin{definition}[Disk Mounting Model]
  \label{def:physics:phyx-mini-0967:phyxminiproblems-problemphyxmini0967-diskmountingmodel}
  \lean{PhyXMiniProblems.ProblemPhyXMini0967.DiskMountingModel}
  How the center of the disk is supported
\end{definition}
\begin{proof}
  This is the declared model definition of Disk Mounting Model.
\end{proof}
\begin{definition}[Disk Scale Regime]
  \label{def:physics:phyx-mini-0967:phyxminiproblems-problemphyxmini0967-diskscaleregime}
  \lean{PhyXMiniProblems.ProblemPhyXMini0967.DiskScaleRegime}
  Scale separation used to treat the disk as a dipole in the axial field
\end{definition}
\begin{proof}
  This is the declared model definition of Disk Scale Regime.
\end{proof}
\begin{definition}[Spin Regime]
  \label{def:physics:phyx-mini-0967:phyxminiproblems-problemphyxmini0967-spinregime}
  \lean{PhyXMiniProblems.ProblemPhyXMini0967.SpinRegime}
  Rapid-spin regime in which the disk axis exhibits gyroscopic precession
\end{definition}
\begin{proof}
  This is the declared model definition of Spin Regime.
\end{proof}
\begin{definition}[Magnetic Interaction Response]
  \label{def:physics:phyx-mini-0967:phyxminiproblems-problemphyxmini0967-magneticinteractionresponse}
  \lean{PhyXMiniProblems.ProblemPhyXMini0967.MagneticInteractionResponse}
  Qualitative response caused by the magnetic interaction
\end{definition}
\begin{proof}
  This is the declared model definition of Magnetic Interaction Response.
\end{proof}
\begin{definition}[Figure Label]
  \label{def:physics:phyx-mini-0967:phyxminiproblems-problemphyxmini0967-figurelabel}
  \lean{PhyXMiniProblems.ProblemPhyXMini0967.FigureLabel}
  Literal parameter and vector labels visible in image 967.png
\end{definition}
\begin{proof}
  This is the declared model definition of Figure Label.
\end{proof}
\begin{definition}[Figure Feature]
  \label{def:physics:phyx-mini-0967:phyxminiproblems-problemphyxmini0967-figurefeature}
  \lean{PhyXMiniProblems.ProblemPhyXMini0967.FigureFeature}
  Physical and graphical objects visible in the supplied raster
\end{definition}
\begin{proof}
  This is the declared model definition of Figure Feature.
\end{proof}
\begin{definition}[Figure Arrow Meaning]
  \label{def:physics:phyx-mini-0967:phyxminiproblems-problemphyxmini0967-figurearrowmeaning}
  \lean{PhyXMiniProblems.ProblemPhyXMini0967.FigureArrowMeaning}
  Physical interpretation of each curved or axial arrow in the figure
\end{definition}
\begin{proof}
  This is the declared model definition of Figure Arrow Meaning.
\end{proof}
\begin{definition}[Rotating Cylinder Disk Figure]
  \label{def:physics:phyx-mini-0967:phyxminiproblems-problemphyxmini0967-rotatingcylinderdiskfigure}
  \lean{PhyXMiniProblems.ProblemPhyXMini0967.RotatingCylinderDiskFigure}
  \uses{def:physics:phyx-mini-0967:phyxminiproblems-problemphyxmini0967-figurelabel, def:physics:phyx-mini-0967:phyxminiproblems-problemphyxmini0967-figurefeature, def:physics:phyx-mini-0967:phyxminiproblems-problemphyxmini0967-figurearrowmeaning}
  Typed transcription of the labels and qualitative incidence information in the 352 × 365 primary raster. It contains no magnetic-field or torque answer readout
\end{definition}
\begin{proof}
  This is the declared model definition of Rotating Cylinder Disk Figure.
\end{proof}
\begin{definition}[Rotating Cylinder Disk Setup]
  \label{def:physics:phyx-mini-0967:phyxminiproblems-problemphyxmini0967-rotatingcylinderdisksetup}
  \lean{PhyXMiniProblems.ProblemPhyXMini0967.RotatingCylinderDiskSetup}
  \uses{def:physics:phyx-mini-0967:phyxminiproblems-problemphyxmini0967-spatialvector, def:physics:phyx-mini-0967:phyxminiproblems-problemphyxmini0967-massmagnitude, def:physics:phyx-mini-0967:phyxminiproblems-problemphyxmini0967-lengthmagnitude, def:physics:phyx-mini-0967:phyxminiproblems-problemphyxmini0967-chargemagnitude, def:physics:phyx-mini-0967:phyxminiproblems-problemphyxmini0967-angularspeedmagnitude, def:physics:phyx-mini-0967:phyxminiproblems-problemphyxmini0967-surfacecurrentdensitymagnitude, def:physics:phyx-mini-0967:phyxminiproblems-problemphyxmini0967-vacuumpermeabilitymagnitude, def:physics:phyx-mini-0967:phyxminiproblems-problemphyxmini0967-magneticfluxdensitymagnitude, def:physics:phyx-mini-0967:phyxminiproblems-problemphyxmini0967-magneticfluxdensityvector, def:physics:phyx-mini-0967:phyxminiproblems-problemphyxmini0967-magneticmomentmagnitude, def:physics:phyx-mini-0967:phyxminiproblems-problemphyxmini0967-magneticmomentvector, def:physics:phyx-mini-0967:phyxminiproblems-problemphyxmini0967-torquemagnitude, def:physics:phyx-mini-0967:phyxminiproblems-problemphyxmini0967-torquevector, def:physics:phyx-mini-0967:phyxminiproblems-problemphyxmini0967-cylindricalshellmodel, def:physics:phyx-mini-0967:phyxminiproblems-problemphyxmini0967-diskmodel, def:physics:phyx-mini-0967:phyxminiproblems-problemphyxmini0967-diskmountingmodel, def:physics:phyx-mini-0967:phyxminiproblems-problemphyxmini0967-diskscaleregime, def:physics:phyx-mini-0967:phyxminiproblems-problemphyxmini0967-spinregime, def:physics:phyx-mini-0967:phyxminiproblems-problemphyxmini0967-magneticinteractionresponse, def:physics:phyx-mini-0967:phyxminiproblems-problemphyxmini0967-rotatingcylinderdiskfigure}
  Independent physical quantities and observables of the cylinder-disk system. The far-from-edge region and cylinder axis are retained as spatial sets. The actual and idealized magnetic fields and torques, along with their physical remainders and error controls, are independent fields constrained only by the governing-law structures below
\end{definition}
\begin{proof}
  This is the declared model definition of Rotating Cylinder Disk Setup.
\end{proof}
\begin{definition}[Matches Rotating Cylinder Disk Scenario]
  \label{def:physics:phyx-mini-0967:phyxminiproblems-problemphyxmini0967-matchesrotatingcylinderdiskscenario}
  \lean{PhyXMiniProblems.ProblemPhyXMini0967.MatchesRotatingCylinderDiskScenario}
  \uses{def:physics:phyx-mini-0967:phyxminiproblems-problemphyxmini0967-rotatingcylinderdisksetup}
  Qualitative apparatus and placement assumptions stated in the problem
\end{definition}
\begin{proof}
  This is the declared model definition of Matches Rotating Cylinder Disk Scenario.
\end{proof}
\begin{definition}[Matches Supplied Rotating Cylinder Disk Figure]
  \label{def:physics:phyx-mini-0967:phyxminiproblems-problemphyxmini0967-matchessuppliedrotatingcylinderdiskfigure}
  \lean{PhyXMiniProblems.ProblemPhyXMini0967.MatchesSuppliedRotatingCylinderDiskFigure}
  \uses{def:physics:phyx-mini-0967:phyxminiproblems-problemphyxmini0967-rotatingcylinderdisksetup}
  Direct qualitative and label evidence from the primary image
\end{definition}
\begin{proof}
  This is the declared model definition of Matches Supplied Rotating Cylinder Disk Figure.
\end{proof}
\begin{definition}[Has Physical Rotating Cylinder Disk Parameters]
  \label{def:physics:phyx-mini-0967:phyxminiproblems-problemphyxmini0967-hasphysicalrotatingcylinderdiskparameters}
  \lean{PhyXMiniProblems.ProblemPhyXMini0967.HasPhysicalRotatingCylinderDiskParameters}
  \uses{def:physics:phyx-mini-0967:phyxminiproblems-problemphyxmini0967-massinkilograms, def:physics:phyx-mini-0967:phyxminiproblems-problemphyxmini0967-lengthinmeters, def:physics:phyx-mini-0967:phyxminiproblems-problemphyxmini0967-chargemagnitudeincoulombs, def:physics:phyx-mini-0967:phyxminiproblems-problemphyxmini0967-angularspeedinradianspersecond, def:physics:phyx-mini-0967:phyxminiproblems-problemphyxmini0967-vacuumpermeabilityinnewtonsperamperesquared, def:physics:phyx-mini-0967:phyxminiproblems-problemphyxmini0967-torquemagnitudeinnewtonmeters, def:physics:phyx-mini-0967:phyxminiproblems-problemphyxmini0967-rotatingcylinderdisksetup}
  Positivity and nondegeneracy of the physical branch shown in the figure
\end{definition}
\begin{proof}
  This is the declared model definition of Has Physical Rotating Cylinder Disk Parameters.
\end{proof}
\begin{definition}[Matches Cylinder Disk Orientation Geometry]
  \label{def:physics:phyx-mini-0967:phyxminiproblems-problemphyxmini0967-matchescylinderdiskorientationgeometry}
  \lean{PhyXMiniProblems.ProblemPhyXMini0967.MatchesCylinderDiskOrientationGeometry}
  \uses{def:physics:phyx-mini-0967:phyxminiproblems-problemphyxmini0967-rotatingcylinderdisksetup}
  Three-dimensional orientation encoded by the primary image and prose. The disk spin axis is the normal to the disk, θ is its unoriented angle to the cylinder axis, and the precession axis is the cylinder axis
\end{definition}
\begin{proof}
  This is the declared model definition of Matches Cylinder Disk Orientation Geometry.
\end{proof}
\begin{definition}[Satisfies Rotating Charged Cylinder Field Laws]
  \label{def:physics:phyx-mini-0967:phyxminiproblems-problemphyxmini0967-satisfiesrotatingchargedcylinderfieldlaws}
  \lean{PhyXMiniProblems.ProblemPhyXMini0967.SatisfiesRotatingChargedCylinderFieldLaws}
  \uses{def:physics:phyx-mini-0967:phyxminiproblems-problemphyxmini0967-lengthinmeters, def:physics:phyx-mini-0967:phyxminiproblems-problemphyxmini0967-chargemagnitudeincoulombs, def:physics:phyx-mini-0967:phyxminiproblems-problemphyxmini0967-angularspeedinradianspersecond, def:physics:phyx-mini-0967:phyxminiproblems-problemphyxmini0967-surfacecurrentdensityinamperespermeter, def:physics:phyx-mini-0967:phyxminiproblems-problemphyxmini0967-vacuumpermeabilityinnewtonsperamperesquared, def:physics:phyx-mini-0967:phyxminiproblems-problemphyxmini0967-magneticfluxdensityinteslas, def:physics:phyx-mini-0967:phyxminiproblems-problemphyxmini0967-magneticfluxdensityvectorinteslas, def:physics:phyx-mini-0967:phyxminiproblems-problemphyxmini0967-rotatingcylinderdisksetup}
  The rotating charged shell acts as an azimuthal surface current K = Q₁ ω₁ / (2 π H). The exactly uniform field B∞ = μ₀ K belongs to the infinite-length idealization; it is not asserted to be the actual finite-shell field merely because the disk is far from the ends
\end{definition}
\begin{proof}
  This is the declared model definition of Satisfies Rotating Charged Cylinder Field Laws.
\end{proof}
\begin{definition}[Satisfies Finite Cylinder Field Approximation]
  \label{def:physics:phyx-mini-0967:phyxminiproblems-problemphyxmini0967-satisfiesfinitecylinderfieldapproximation}
  \lean{PhyXMiniProblems.ProblemPhyXMini0967.SatisfiesFiniteCylinderFieldApproximation}
  \uses{def:physics:phyx-mini-0967:phyxminiproblems-problemphyxmini0967-magneticfluxdensityinteslas, def:physics:phyx-mini-0967:phyxminiproblems-problemphyxmini0967-magneticfluxdensityvectorinteslas, def:physics:phyx-mini-0967:phyxminiproblems-problemphyxmini0967-rotatingcylinderdisksetup}
  Finite-cylinder approximation contract. The actual field at the disk center is decomposed into the ideal infinite-shell field and a stored vector remainder. Its norm is controlled by a dimensionless relative tolerance. This is an explicit finite-parameter error statement, not an exact field law selected by a qualitative regime tag
\end{definition}
\begin{proof}
  This is the declared model definition of Satisfies Finite Cylinder Field Approximation.
\end{proof}
\begin{definition}[Satisfies Spinning Charged Disk Magnetic Moment Law]
  \label{def:physics:phyx-mini-0967:phyxminiproblems-problemphyxmini0967-satisfiesspinningchargeddiskmagneticmomentlaw}
  \lean{PhyXMiniProblems.ProblemPhyXMini0967.SatisfiesSpinningChargedDiskMagneticMomentLaw}
  \uses{def:physics:phyx-mini-0967:phyxminiproblems-problemphyxmini0967-lengthinmeters, def:physics:phyx-mini-0967:phyxminiproblems-problemphyxmini0967-chargemagnitudeincoulombs, def:physics:phyx-mini-0967:phyxminiproblems-problemphyxmini0967-angularspeedinradianspersecond, def:physics:phyx-mini-0967:phyxminiproblems-problemphyxmini0967-magneticmomentmagnitudeinamperesquaremeters, def:physics:phyx-mini-0967:phyxminiproblems-problemphyxmini0967-magneticmomentvectorinamperesquaremeters, def:physics:phyx-mini-0967:phyxminiproblems-problemphyxmini0967-rotatingcylinderdisksetup}
  The supplied magnetic moment of the uniformly charged thin disk is μ = Q₂ ω₂ R₂² / 4, directed along the disk's spin axis
\end{definition}
\begin{proof}
  This is the declared model definition of Satisfies Spinning Charged Disk Magnetic Moment Law.
\end{proof}
\begin{definition}[Satisfies Magnetic Dipole Torque Law]
  \label{def:physics:phyx-mini-0967:phyxminiproblems-problemphyxmini0967-satisfiesmagneticdipoletorquelaw}
  \lean{PhyXMiniProblems.ProblemPhyXMini0967.SatisfiesMagneticDipoleTorqueLaw}
  \uses{def:physics:phyx-mini-0967:phyxminiproblems-problemphyxmini0967-magneticfluxdensityvectorinteslas, def:physics:phyx-mini-0967:phyxminiproblems-problemphyxmini0967-magneticmomentvectorinamperesquaremeters, def:physics:phyx-mini-0967:phyxminiproblems-problemphyxmini0967-torquemagnitudeinnewtonmeters, def:physics:phyx-mini-0967:phyxminiproblems-problemphyxmini0967-torquevectorinnewtonmeters, def:physics:phyx-mini-0967:phyxminiproblems-problemphyxmini0967-spatialcross, def:physics:phyx-mini-0967:phyxminiproblems-problemphyxmini0967-rotatingcylinderdisksetup}
  In the uniform infinite-shell idealization, the magnetic torque is τ∞ = μ × B∞; its separately stored magnitude is the Euclidean norm of that ideal vector observable
\end{definition}
\begin{proof}
  This is the declared model definition of Satisfies Magnetic Dipole Torque Law.
\end{proof}
\begin{definition}[Satisfies Finite Apparatus Torque Approximation]
  \label{def:physics:phyx-mini-0967:phyxminiproblems-problemphyxmini0967-satisfiesfiniteapparatustorqueapproximation}
  \lean{PhyXMiniProblems.ProblemPhyXMini0967.SatisfiesFiniteApparatusTorqueApproximation}
  \uses{def:physics:phyx-mini-0967:phyxminiproblems-problemphyxmini0967-torquemagnitudeinnewtonmeters, def:physics:phyx-mini-0967:phyxminiproblems-problemphyxmini0967-torquevectorinnewtonmeters, def:physics:phyx-mini-0967:phyxminiproblems-problemphyxmini0967-rotatingcylinderdisksetup}
  The actual finite-apparatus torque is the ideal uniform-field dipole torque plus a physical vector remainder. The remainder bound accounts for the finite cylinder, residual field variation across the small disk, and any higher multipole contribution suppressed by the scale separation. It does not assume the answer formula
\end{definition}
\begin{proof}
  This is the declared model definition of Satisfies Finite Apparatus Torque Approximation.
\end{proof}
\begin{lemma}[rotating Cylinder Ideal Interior Field Magnitude Formula]
  \label{lem:physics:phyx-mini-0967:phyxminiproblems-problemphyxmini0967-rotatingcylinderidealinteriorfieldmagnitudeformula}
  \lean{PhyXMiniProblems.ProblemPhyXMini0967.rotatingCylinderIdealInteriorFieldMagnitudeFormula}
  \uses{def:physics:phyx-mini-0967:phyxminiproblems-problemphyxmini0967-lengthinmeters, def:physics:phyx-mini-0967:phyxminiproblems-problemphyxmini0967-chargemagnitudeincoulombs, def:physics:phyx-mini-0967:phyxminiproblems-problemphyxmini0967-angularspeedinradianspersecond, def:physics:phyx-mini-0967:phyxminiproblems-problemphyxmini0967-vacuumpermeabilityinnewtonsperamperesquared, def:physics:phyx-mini-0967:phyxminiproblems-problemphyxmini0967-magneticfluxdensityinteslas, def:physics:phyx-mini-0967:phyxminiproblems-problemphyxmini0967-rotatingcylinderdisksetup, def:physics:phyx-mini-0967:phyxminiproblems-problemphyxmini0967-satisfiesrotatingchargedcylinderfieldlaws}
  The two rotating-shell laws imply the ideal infinite-shell axial field
\end{lemma}
\begin{proof}
  The conclusion follows from the governing relations and figure readouts named in the statement of rotating Cylinder Ideal Interior Field Magnitude Formula.
\end{proof}
\begin{lemma}[ideal Magnetic Dipole Torque Magnitude Formula]
  \label{lem:physics:phyx-mini-0967:phyxminiproblems-problemphyxmini0967-idealmagneticdipoletorquemagnitudeformula}
  \lean{PhyXMiniProblems.ProblemPhyXMini0967.idealMagneticDipoleTorqueMagnitudeFormula}
  \uses{def:physics:phyx-mini-0967:phyxminiproblems-problemphyxmini0967-magneticfluxdensityinteslas, def:physics:phyx-mini-0967:phyxminiproblems-problemphyxmini0967-magneticmomentmagnitudeinamperesquaremeters, def:physics:phyx-mini-0967:phyxminiproblems-problemphyxmini0967-torquemagnitudeinnewtonmeters, def:physics:phyx-mini-0967:phyxminiproblems-problemphyxmini0967-rotatingcylinderdisksetup, def:physics:phyx-mini-0967:phyxminiproblems-problemphyxmini0967-matchescylinderdiskorientationgeometry, def:physics:phyx-mini-0967:phyxminiproblems-problemphyxmini0967-satisfiesrotatingchargedcylinderfieldlaws, def:physics:phyx-mini-0967:phyxminiproblems-problemphyxmini0967-satisfiesspinningchargeddiskmagneticmomentlaw, def:physics:phyx-mini-0967:phyxminiproblems-problemphyxmini0967-satisfiesmagneticdipoletorquelaw}
  Before substituting either source formula, the ideal dipole law and angle geometry give |τ∞| = μ B∞ sin θ
\end{lemma}
\begin{proof}
  The conclusion follows from the governing relations and figure readouts named in the statement of ideal Magnetic Dipole Torque Magnitude Formula.
\end{proof}
\begin{lemma}[finite Apparatus Torque Magnitude Error Bound]
  \label{lem:physics:phyx-mini-0967:phyxminiproblems-problemphyxmini0967-finiteapparatustorquemagnitudeerrorbound}
  \lean{PhyXMiniProblems.ProblemPhyXMini0967.finiteApparatusTorqueMagnitudeErrorBound}
  \uses{def:physics:phyx-mini-0967:phyxminiproblems-problemphyxmini0967-torquemagnitudeinnewtonmeters, def:physics:phyx-mini-0967:phyxminiproblems-problemphyxmini0967-rotatingcylinderdisksetup, def:physics:phyx-mini-0967:phyxminiproblems-problemphyxmini0967-satisfiesfiniteapparatustorqueapproximation}
  The reverse triangle inequality transports the vector-remainder contract to an absolute error bound for torque magnitudes. This is the quantitative finite-apparatus conclusion that replaces the former global exact equality
\end{lemma}
\begin{proof}
  The conclusion follows from the governing relations and figure readouts named in the statement of finite Apparatus Torque Magnitude Error Bound.
\end{proof}
\begin{definition}[Answer Choice]
  \label{def:physics:phyx-mini-0967:phyxminiproblems-problemphyxmini0967-answerchoice}
  \lean{PhyXMiniProblems.ProblemPhyXMini0967.AnswerChoice}
  Labels of the four torque-magnitude choices printed with the problem
\end{definition}
\begin{proof}
  This is the declared model definition of Answer Choice.
\end{proof}
\begin{definition}[displayed Torque Magnitude In Newton Meters]
  \label{def:physics:phyx-mini-0967:phyxminiproblems-problemphyxmini0967-displayedtorquemagnitudeinnewtonmeters}
  \lean{PhyXMiniProblems.ProblemPhyXMini0967.displayedTorqueMagnitudeInNewtonMeters}
  \uses{def:physics:phyx-mini-0967:phyxminiproblems-problemphyxmini0967-lengthinmeters, def:physics:phyx-mini-0967:phyxminiproblems-problemphyxmini0967-chargemagnitudeincoulombs, def:physics:phyx-mini-0967:phyxminiproblems-problemphyxmini0967-angularspeedinradianspersecond, def:physics:phyx-mini-0967:phyxminiproblems-problemphyxmini0967-vacuumpermeabilityinnewtonsperamperesquared, def:physics:phyx-mini-0967:phyxminiproblems-problemphyxmini0967-rotatingcylinderdisksetup, def:physics:phyx-mini-0967:phyxminiproblems-problemphyxmini0967-answerchoice}
  Scalar coherent-SI expressions displayed beside each answer label. Choice A is transcribed literally even though it omits one power of R₂
\end{definition}
\begin{proof}
  This is the declared model definition of displayed Torque Magnitude In Newton Meters.
\end{proof}
\begin{definition}[recorded Dataset Answer]
  \label{def:physics:phyx-mini-0967:phyxminiproblems-problemphyxmini0967-recordeddatasetanswer}
  \lean{PhyXMiniProblems.ProblemPhyXMini0967.recordedDatasetAnswer}
  \uses{def:physics:phyx-mini-0967:phyxminiproblems-problemphyxmini0967-answerchoice}
  The answer label recorded in the source dataset
\end{definition}
\begin{proof}
  This is the declared model definition of recorded Dataset Answer.
\end{proof}
% --- Archon declaration topology end ---
% --- Archon physics formalization source end ---
